% LaTeX Template for short student reports.
% Citations should be in bibtex format and go in references.bib
\documentclass[a4paper, 11pt]{article}
\usepackage[top=3cm, bottom=3cm, left = 2cm, right = 2cm]{geometry} 
\geometry{a4paper} 
\usepackage[utf8]{inputenc}
\usepackage{textcomp}
\usepackage{graphicx} 
\usepackage{amsmath,amssymb}  
\usepackage{bm}  
\usepackage[pdftex,bookmarks,colorlinks,breaklinks]{hyperref}  
%\hypersetup{linkcolor=black,citecolor=black,filecolor=black,urlcolor=black} % black links, for printed output
\usepackage{memhfixc} 
\usepackage{pdfsync}  
\usepackage{fancyhdr}
\pagestyle{fancy}

\title{A (possibly) Interesting set of SAT problems}
\author{Timothy McMahon}
%\date{}

\begin{document}
%\maketitle
%\tableofcontents

\section{Introduction}

This is a list of problems that should be solveable using SAT.  Hopefully they are of interest and lead to some interesting decompositions for dagster.

%\pagebreak

\section{Domatic Partitions of the Boolean Cube}
 
The $n$-dimensional Boolean cube is the simple graph $Q_n$, whose vertices $V$ are the sequences of $0$s and $1$s of length $n$.  There is an edge between two vertices if the Hamming distance between them is 1 (i.e. if they differ in exactly one bit). 

\begin{figure}[h!]
    \centering
     \includegraphics[width=0.5\textwidth]{{boolean_cube.drawio.pdf}}
    \caption{The first few Boolean cubes}
    \end{figure}


A subset $S$ of the vertices is \textbf{dominating} if every vertex not in $S$ is adjacent to some vertex in $S$.
A partition of the vertices of $Q_n$ into sets $S_1, S_2, \ldots , S_k$ is a \textbf{domatic} partition if each set is dominating.  
 
We define $\Tau(n)$ to be the largest $k$ such that $Q_n$ has a domatic partition into $k$ sets.  

For example, $\Tau(3) = 4$, as the vertices of the $3$-dimensional cube can be partitioned into 4 dominating sets, $\{000, 111\}, \{001, 110\}, \{010, 101\}, \{011, 100\}, but not into 5.$
 
The known values of $\Tau(n)$ are below.

\begin{center}
    \begin{tabular}{ | c | c | }
    \hline
     2 & 2 \\ 
     3 & 4 \\ 
     4 & 4 \\ 
     5 & 4 \\ 
     6 & 5 \\ 
     7 & 8 \\ 
     8 & 8 \\ 
     9 & 8 \\ 
     10 & (9 or 10) ? \\ 
    \end{tabular}
    \end{center}



This is the OEIS sequence A157887.  
There is a comment on the OEIS page that says that the domatic number for 10 must be 8 or 9, but I didn't follow the argument, so maybe this has been solved? The initial place I read it claimed that a partition for 9 had been found, but they couldn't rule out 10.  I'll try to track down some references....
 
 
Graceful labelings of graphs
 
Label the vertices of a tree from 0 to n-1.   This induces a labelling on the edges : the edge label is the difference between its endpoints.
A labelling is graceful if the edges are all distinct.
 
Conjecture (Ringel-Kotzig): All trees are graceful (i.e. they admit a graceful labelling).
 
Further, all trees are bipartite, so colour all of the vertices in one part white and the other black.  A graceful colouring is bipartite if for every edge, the black endpoint has a higher number than the white endpoint.
 
Conjecture (Pulleyblank):  All trees admit a bipartite graceful colouring.
 
The specific 81 vertex graph that they hadn't found a labelling (but potentially has been found now, the reference was from 2014) is
 
v1: (1, 2, 6)
v2: (3, 4, 0)
v3: (2, 0, 6)
v4: (4, 3, 5)
v5: (0, 2, 2, 3)
v6: ()
v7: (0, 1, 1, 4, 4)
 
 
The first 7 vertices (v1 - v7) are in a path (i.e  v1 ---- v2 ---- v3 ---- v4 --- ...  v7), and the notation v1: (1, 2, 6) means v1 has 3 children, which in turn have 1, 2 and 6 children respectively.  This type of graph is called a 'lobster' graph.

\pagebreak

\section{System Design}

My system worked as follows \ldots

\pagebreak

\section{Evaluation}

We did some experiments \ldots

\pagebreak

\section{Conclusions and Future Work}

From our experiments we can conclude that \ldots

\bibliographystyle{abbrv}
% \bibliography{references}  % need to put bibtex references in references.bib 
\end{document}